\section{Conclusioni}
In questo lavoro sono stati implementati e
confrontati diversi modelli CNN per la super-resolution ispirati
a SRCNN\cite{dong2014srcnn} e VDSR\cite{kim2015accurate}.
 Nonostante i miglioramenti progressivi introdotti,
 le prestazioni non hanno superato sistematicamente l'interpolazione bicubica.
 Il risultato migliore è stato ottenuto con la loss L1,
 che ha permesso di superare Bicubic in un caso del test set con un
 miglioramento percettivo visibile sui dettagli.
 La difficoltà nel raggiungere prestazioni superiori è probabilmente
 attribuibile all'utilizzo di un dataset di training diverso rispetto ai paper di riferimento,
 all'assenza di data augmentation e a una profondità della rete ridotta.
 Il fatto che la loss si stabilizzasse dopo circa cento epoche suggerisce
 inoltre che il modello raggiungeva i limiti della propria capacità rappresentativa,
 indipendentemente dalla durata dell'addestramento.